% !TEX root = ../Main.tex
\chapter{精选题}

本章汇集两道综合题，综合运用前几章方法——计数原理、对立事件、独立重复试验、二项分布。

\section{多选题的概率（多选）}

\ex[多选题的随机猜]{
    某次数学考试的一道\textbf{多项选择题}，要求是："在每小题给出的四个选项中，\textbf{全部选对}的得 $5$ 分，\textbf{部分选对}的得 $3$ 分，\textbf{有选错}的得 $0$ 分。"已知某选择题的正确答案是 $CD$，且甲、乙、丙、丁四位同学都不会做，下列表述正确的是哪些？
    \begin{enumerate}[label=(\Alph*)]
        \item 甲仅随机选 $1$ 个选项，能得 $3$ 分的概率是 $\tfrac{1}{2}$；
        \item 乙仅随机选 $2$ 个选项，能得 $5$ 分的概率是 $\tfrac{1}{6}$；
        \item 丙随机选择选项（从非空子集中等概率选），能得分的概率是 $\tfrac{1}{5}$；
        \item 丁随机至少选择 $2$ 个选项，能得分的概率是 $\tfrac{1}{10}$。
    \end{enumerate}
}

\sol{
    正确答案集 $\{C, D\}$。选项共 $4$ 个，正确的 $2$ 个、错误的 $2$ 个。"得分" $\iff$ 所选都在 $\{C, D\}$ 中且非空。

    \textbf{(A)} 甲随机选 $1$ 个，样本空间 $\{A, B, C, D\}$ 四种等可能。选中 $C$ 或 $D$ 即得 $3$ 分（部分选对）——$2$ 种。
    \[
    P_A = \frac{2}{4} = \frac{1}{2}.\quad \checkmark
    \]

    \textbf{(B)} 乙随机选 $2$ 个，样本空间 $C_4^2 = 6$ 种。得 $5$ 分 $\iff$ 选中 $\{C, D\}$——$1$ 种。
    \[
    P_B = \frac{1}{6}.\quad \checkmark
    \]

    \textbf{(C)} 丙"随机选"——按通常解读，从所有\textbf{非空子集}中等概率选一个，共 $2^4 - 1 = 15$ 种。"能得分" $\iff$ 子集 $\subseteq \{C, D\}$ 且非空——即 $\{C\}, \{D\}, \{C, D\}$ 共 $3$ 种。
    \[
    P_C = \frac{3}{15} = \frac{1}{5}.\quad \checkmark
    \]

    \textbf{(D)} 丁随机"至少选 $2$ 个"：选 $2, 3, 4$ 个的子集共 $C_4^2 + C_4^3 + C_4^4 = 6 + 4 + 1 = 11$ 种。"能得分" $\iff$ 子集 $\subseteq \{C, D\}$。
    \begin{itemize}
        \item 选 $2$ 个：$\{C, D\}$——$1$ 种（得 $5$ 分）；
        \item 选 $3$ 个：需 $3$ 个都在 $\{C, D\}$ 内，但 $\{C, D\}$ 只有 $2$ 个元素——$0$ 种；
        \item 选 $4$ 个：同理——$0$ 种。
    \end{itemize}
    合计 $1$ 种。
    \[
    P_D = \frac{1}{11} \ne \frac{1}{10}.\quad \times
    \]

    \textbf{正确选项：(A)(B)(C)}。
}

\section{NBA 七场四胜制}

\ex[恰好 $5$ 场决出冠军]{
    NBA 总决赛采用 $7$ 场 $4$ 胜制（谁先赢 $4$ 场谁夺冠）。$2018$ 年总决赛两支球队分别为勇士和骑士，假设每场比赛勇士获胜的概率为 $0.6$、骑士获胜的概率为 $0.4$，且每场比赛结果相互独立。求\textbf{恰好 $5$ 场}比赛决出总冠军的概率。
}

\sol{
    "恰好 $5$ 场结束" $\iff$ 第 $5$ 场某方\textbf{刚好取得第 $4$ 胜}，且此前 $4$ 场该方胜 $3$ 场、另一方胜 $1$ 场。

    \textbf{情形 1：勇士以 $4:1$ 夺冠。} 前 $4$ 场勇士胜 $3$ 场（$C_4^3$ 种排列），第 $5$ 场勇士胜：
    \[
    P_1 = C_4^3 \cdot 0.6^3 \cdot 0.4 \cdot 0.6 = 4 \cdot 0.6^4 \cdot 0.4.
    \]

    \textbf{情形 2：骑士以 $4:1$ 夺冠。} 前 $4$ 场骑士胜 $3$ 场，第 $5$ 场骑士胜：
    \[
    P_2 = C_4^3 \cdot 0.4^3 \cdot 0.6 \cdot 0.4 = 4 \cdot 0.4^4 \cdot 0.6.
    \]

    \textbf{两情形互斥，求和：}
    \[
    P = 4 \cdot 0.6^4 \cdot 0.4 + 4 \cdot 0.4^4 \cdot 0.6 = 4 \cdot 0.6 \cdot 0.4 \cdot (0.6^3 + 0.4^3).
    \]
    计算：$0.6^3 = 0.216$，$0.4^3 = 0.064$，和为 $0.28$；$4 \cdot 0.6 \cdot 0.4 = 0.96$。
    \[
    P = 0.96 \cdot 0.28 = 0.2688 = \frac{168}{625}.
    \]
}

\rmkb{
    \textbf{"恰好 $n$ 场结束"类题目的通用套路}：
    \begin{enumerate}
        \item 先固定\textbf{最后一场}——是某方取得"决胜手"（第 $w$ 胜）的那一场；
        \item 前 $n - 1$ 场中该方胜 $w - 1$ 场、另一方胜 $n - w$ 场——\textbf{组合数}给出排列方式；
        \item 总冠军两队都可能——\textbf{按互斥求和}。
    \end{enumerate}

    \textbf{常见错误}：
    \begin{itemize}
        \item 忘记\textbf{最后一场必须是决胜方赢}——这是"恰好"的关键约束；
        \item 忘记乘\textbf{组合数}——前 $n-1$ 场的胜场排列方式不唯一；
        \item 漏掉\textbf{另一队夺冠}的情形——两种都是"恰好 $n$ 场"。
    \end{itemize}
}

\section{综合题的破题思路}

\rmkb{
    概率综合题的三步诊断：
    \begin{enumerate}
        \item \textbf{辨识题型}：古典概型、二项分布、条件概率、方差……；
        \item \textbf{选择工具}：计数原理、对立事件、乘法公式、全概率、方差计算式；
        \item \textbf{构造方程 / 不等式}：把题目所求用上述工具表达、求解。
    \end{enumerate}

    \textbf{常见条件反射（熟练度）}：
    \begin{itemize}
        \item 见"至少" → 想\textbf{对立事件}；
        \item 见"恰好 $k$ 次成功" → 想\textbf{二项分布 + 组合数}；
        \item 见"独立" → 想\textbf{乘积公式}；
        \item 见"不放回抽取" → 想\textbf{组合 / 排列样本空间}；
        \item 见"已知 A 发生，问 B 的概率" → 想\textbf{条件概率}；
        \item 见"平均数 + 方差" → 想 $\overline{x^2} - \bar x^2$ 或 $\sum (x_i - \bar x) = 0$。
    \end{itemize}
    这些反射是大量做题积累的——概率题的"难"往往不在计算，而在\textbf{识别出正确的工具}。
}
